\section{Overflow}

Properti \code{overflow} mengontrol perilaku konten yang melebihi dimensi kotak elemen. Situasi ini sering terjadi ketika elemen memiliki \code{height} atau \code{width} yang tetap sementara kontennya lebih besar. Memahami overflow penting untuk mencegah layout yang rusak dan menyediakan pengalaman pengguna yang baik pada konten dinamis \cite{mdnCss}.

\begin{table}[htbp]
\centering
\begin{tabular}{|P{3cm}|P{9cm}|}
\hline
\textbf{Nilai} & \textbf{Perilaku} \\
\hline
\code{visible} (default) & Konten meluap keluar kotak; dapat menutupi elemen lain \\
\hline
\code{hidden} & Konten yang meluap dipotong; tidak ada scrollbar \\
\hline
\code{scroll} & Scrollbar selalu ditampilkan, meskipun konten tidak meluap \\
\hline
\code{auto} & Scrollbar hanya muncul jika konten meluap \\
\hline
\code{clip} & Seperti hidden, tetapi melarang scroll programatis \\
\hline
\end{tabular}
\caption{Nilai properti overflow dan perilakunya}
\end{table}

Properti \code{overflow-x} dan \code{overflow-y} mengontrol overflow secara terpisah untuk sumbu horizontal dan vertikal. Kombinasi \code{overflow-x: hidden; overflow-y: auto;} sering digunakan pada sidebar atau panel navigasi. Teknik \code{text-overflow: ellipsis} dengan \code{overflow: hidden; white-space: nowrap;} memotong teks panjang dengan tanda elipsis (``...'') \cite{webdevCss}.

\begin{konsep}
\textbf{Overflow dan Block Formatting Context.} Elemen dengan \code{overflow} selain \code{visible} menciptakan \textbf{Block Formatting Context} (BFC) baru. Ini berarti elemen tersebut akan memuat float di dalamnya dan margin-nya tidak akan collapse dengan elemen luar. Properti \code{display: flow-root} menciptakan BFC tanpa efek samping overflow \cite{mdnCss}.
\end{konsep}

